\chapter{核心流水线、乱序执行与分支预测}

现代 CPU 核心不是按源码顺序一条一条机械执行。为了提高吞吐，核心会把指令拆成微操作，放入队列，分析依赖，尽量让互不依赖的操作并行执行。理解这一层，是理解单线程性能、SIMD 前的标量优化和并发热路径的基础。

\section{流水线不是免费并行}

流水线把取指、译码、执行、访存和提交分成阶段。理想情况下，每个周期都有新指令进入，每个阶段都在工作。但真实程序有数据依赖、控制依赖和资源冲突。后一条指令需要前一条指令结果时，执行端口即使空闲也不能使用；分支目标未知时，前端可能取错指令；访存 miss 时，后端会等待数据。

流水线的关键不是“越深越快”，而是吞吐和延迟的权衡。深流水线可以提高频率，却放大分支预测失败的代价。并发程序中的短临界区也会受流水线影响：一个看似只有几条指令的锁操作，背后可能包含原子读改写、cache line 独占获取和内存序约束。

\section{前端：取指、译码和微操作}

核心前端负责把指令流喂给后端。它要从 instruction cache 取指，预测下一条指令地址，译码复杂指令，把它们转成内部微操作，并放入队列。前端跟不上时，后端即使有空闲执行端口也没有工作可做。大函数、过多分支、间接跳转、指令 cache miss 和译码压力都可能让前端成为瓶颈。

x86-64 指令长度可变，译码比定长指令集复杂。现代核心通常会使用 micro-op cache 或类似结构缓存译码结果，让热循环不用每次重新译码。这个机制解释了为什么热循环体大小、函数内联和代码布局会影响性能。内联可以减少调用开销，但过度内联会膨胀代码，伤害 instruction cache 和前端供给。

第一册已经教读者读汇编。第二册读汇编时要多看一个维度：这段机器码对前端是否友好。短小热循环、稳定分支、连续代码布局和可预测调用目标，通常更容易让前端持续供给后端。反过来，大量模板实例化、过度内联、异常路径混在热路径、虚调用目标多变、解释器式大 switch，都可能让前端成本超过源码直觉。

\begin{deepdive}
前端瓶颈的一个典型现象是：数据已经在缓存中，算术也不复杂，但 IPC 仍然低，并且改变数据布局帮助不大。这时应检查代码体积、分支目标、间接跳转和译码压力，而不是继续只改数组排列。
\end{deepdive}

\section{后端：执行端口和调度}

后端接收微操作后，要等待操作数准备好，并把它们发射到合适执行端口。整数 ALU、浮点乘加、加载、存储、分支和地址生成都有各自资源。一个循环如果每轮需要太多 load，可能受加载端口限制；如果每轮地址计算复杂，可能受地址生成单元限制；如果依赖链很长，端口空着也不能发射。

因此分析单核性能时，要区分“端口吞吐不够”和“依赖导致发不出去”。多累加器能改善后者；SIMD 能提高每条指令处理的数据量；简化地址表达式和连续访问能减轻加载和地址生成压力。不同 CPU 微架构端口配置不同，严肃优化必须结合具体机器。

后端不是一个抽象黑箱。一个 load 通常需要地址生成、TLB 查询、缓存访问和可能的 miss 处理；一个 store 通常拆成地址和数据路径，并进入 store buffer；一个浮点 FMA 会占用特定向量执行资源；一个整数分支会占用分支执行资源。循环的性能取决于这些微操作怎样分布到端口上，也取决于它们之间是否互相等待。

这解释了为什么源码层面的“操作次数”常常误导。两个循环都执行同样数量的加法，一个可能因为单累加器依赖链而慢，另一个可能因为多累加器让后端并行发射而快。两个循环都做同样数量的 load，一个可能连续访问被预取器支持，另一个可能指针追踪导致每次 load 等上一次 load 的地址。

\section{乱序执行的边界}

乱序执行允许核心在不改变单线程可观察结果的前提下，提前执行已经准备好的指令。寄存器重命名消除了很多假依赖，保留站和重排序缓冲区让多个指令同时等待资源。这样做可以隐藏一部分延迟，例如在一次乘法等待时执行另一个独立加法。

但是乱序执行不能越过所有边界。真实数据依赖不能消除，cache miss 太多时可并行等待的 miss 数量有限，内存屏障会限制重排，原子操作会带来更强的顺序约束。写高性能代码时，增加独立累加器、减少循环携带依赖、让数据访问可预测，都是在帮助乱序核心找到可调度空间。

\section{重排序缓冲区和执行窗口}

乱序执行不是无限的。核心只能观察一个有限窗口内的指令。重排序缓冲区保存尚未提交的指令，load buffer 保存未完成加载，store buffer 保存尚未对外可见的存储。如果窗口被长延迟 load 占住，后续可执行工作又不够，核心会停顿。

链表遍历就是典型例子。每次读取下一个指针之前，必须等当前节点加载完成。乱序核心很难提前知道下一步地址，所以内存级并行度很低。相比之下，遍历多个独立数组或同时处理多条链，可以让多个 load 同时飞行，更容易隐藏延迟。

执行窗口可以被理解为硬件能同时“看见”的未来。窗口越大，越有机会在一个长延迟操作等待时找到其他独立工作。但如果代码中充满循环携带依赖，窗口再大也找不到可提前执行的操作；如果 cache miss 太多，load buffer 被占满，后续加载也发不出去；如果分支预测失败，窗口里的错误路径工作会被丢弃。

软件能做的事情，是把独立工作显式放进窗口。多累加器、分块处理、批量处理多个请求、一次处理多个链表游标、把检查从热循环中外提，都属于给乱序核心提供更多可调度空间。相反，频繁调用不可内联函数、隐藏别名、每步都依赖共享原子结果，会压缩这个空间。

\section{Load、Store 和内存消歧}

内存操作比寄存器操作复杂。核心要判断一个 load 是否可能读取前面尚未完成的 store。如果地址关系不明确，硬件会做内存消歧预测。预测错误时，需要回滚并重新执行相关操作。高性能循环中，减少指针别名、使用清晰的数组边界、避免复杂交叉写读，可以帮助硬件和编译器更大胆地调度。

store buffer 也影响并发语义。普通 store 可以先进入 store buffer，稍后再对其他核心可见。内存屏障、原子操作和锁释放会约束这种行为。C++ 内存模型的 acquire/release 最终要落到这类硬件机制上。

Load/store queue 还是很多微妙性能问题的来源。若一个循环中既写数组又读另一个可能别名的数组，编译器和硬件都必须保守处理读写顺序。若 store 地址很晚才算出来，后面的 load 就难以确认是否可以提前。若 store buffer 堆积，后续 store 可能被阻塞，锁释放和原子操作也可能放大这个问题。

并发程序中，store buffer 让“本线程已经执行了写”不等于“其他线程已经能按你想象的顺序看到写”。这不是说硬件不可靠，而是说可见性需要同步语义。后面讲 release/acquire 时，会把这条硬件直觉变成 C++ 层面的 happens-before 关系。

\section{分支预测和控制流}

分支预测让核心在条件结果出来之前先猜路径。预测正确时，流水线保持饱满；预测失败时，错误路径上的工作要被丢弃。随机分支、数据相关分支和复杂间接跳转都会破坏预测。很多高性能循环会把条件变成掩码、查表或分段处理，就是为了减少不可预测控制流。

并发程序也受分支预测影响。锁的 fast path 通常期望“不竞争”，如果竞争比例上升，分支行为和 cache 行为都会改变。无锁结构的 CAS 循环在低竞争时很短，在高竞争时会反复失败，控制流和内存系统同时恶化。

\begin{deepdive}
判断一段代码是不是受分支预测影响，可以比较不同输入分布下的 cycles、branches 和 branch-misses。如果随机数据显著慢于有序数据，而内存访问量没有明显变化，分支预测通常就是主要嫌疑。
\end{deepdive}

\section{SMT 和共享核心资源}

SMT 让一个物理核心暴露多个逻辑 CPU。它可以在一个线程等待内存或前端供给不足时，让另一个线程使用空闲资源。但 SMT 不是免费翻倍。两个逻辑线程会共享前端、执行端口、缓存、TLB 和乱序资源的一部分。如果两个线程都在做同样高强度向量计算，它们可能互相抢执行资源；如果一个线程经常等内存，另一个线程可能填补空档。

因此 benchmark 中不能只写“8 个 CPU”。要区分物理核心和 SMT 线程。对于 CPU-bound 数值内核，先扩展到物理核心往往比立刻使用所有逻辑 CPU 更稳定；对于延迟隐藏明显的工作负载，SMT 可能提高吞吐。结论必须来自本机测量，而不是固定经验。

\section{微架构和 ISA 的边界}

第一册强调 ISA 和 ABI：同一条 \instruction{add} 在 x86-64 语义上做什么，函数参数怎样传递。第二册强调微架构：同一个 ISA 指令在不同 CPU 上可能拆成不同微操作，使用不同端口，拥有不同延迟和吞吐。ISA 给出程序可见语义，微架构决定实现成本。

这条边界非常重要。写教材时不能把某个 CPU 的端口配置说成 x86-64 的永久性质；写性能报告时也不能只说“x86 上这个指令很快”。更准确的说法是：在某个 CPU 型号、某个编译选项、某个输入分布下，观察到某个指令序列的瓶颈更像前端、端口、依赖、内存或分支。

\section{典型性能症状}

前端瓶颈常表现为 IPC 低、branch miss 或 instruction cache 相关事件高，代码布局和热循环大小值得检查。后端执行瓶颈常表现为数据在缓存中、分支也稳定，但 IPC 仍受端口或依赖限制，适合检查 SIMD、多累加器和指令混合。内存瓶颈常表现为 cache miss、TLB miss 或带宽接近平台上限。同步瓶颈则可能表现为 IPC 低、cache line 迁移、context switch、futex 等待或原子热点。

这些症状不会自动给答案，但能排除方向。不要在内存带宽瓶颈上执着减少几条整数指令，也不要在锁竞争瓶颈上只调整循环展开。性能工程的价值在于把“慢”拆成可以证伪的假设。

\section{从硬件回到代码}

核心流水线给软件三个直接启发。第一，减少长依赖链，例如把单个累加器拆成多个局部累加器。第二，减少不可预测分支，让热路径尽量直。第三，减少会强制顺序的操作，例如不必要的原子和屏障。后续章节讨论 SIMD、锁、无锁队列和并行归约时，都会回到这些原则。

\begin{labbox}
实验：写两个求和版本。第一个使用单累加器，第二个使用四个独立累加器。用 \cmd{perf stat} 比较 cycles、instructions 和 IPC。再构造一个带随机分支的版本，观察 branch-misses 对时间的影响。
\end{labbox}
